\documentclass[11pt, letterpaper]{article}

% ----- PAQUETES -----
\usepackage[utf8]{inputenc}
\usepackage[spanish]{babel}
\usepackage{graphicx}
\usepackage[margin=2.3cm]{geometry}
\usepackage{xcolor}
\usepackage{titlesec}
\usepackage{booktabs}
\usepackage{enumitem}
\usepackage{xurl}
\usepackage{hyperref}
\usepackage{fancyhdr}
\usepackage{amsmath}
\usepackage{amssymb}

% ----- COLORES -----
\definecolor{UnalGreen}{RGB}{0,153,51}
\definecolor{DarkGray}{RGB}{60,60,60}

% ----- SECCIONES -----
\titleformat{\section}
{\normalfont\large\bfseries\color{UnalGreen}}
{\thesection}{1em}{}[{\titlerule[0.8pt]}]

\titleformat{\subsection}
{\normalfont\normalsize\bfseries\color{DarkGray}}
{\thesubsection}{1em}{}

\titlespacing*{\section}{0pt}{8pt}{4pt}
\titlespacing*{\subsection}{0pt}{6pt}{2pt}

% ----- ENLACES -----
\hypersetup{
	colorlinks=true,
	linkcolor=UnalGreen,
	urlcolor=UnalGreen
}

% ----- ENCABEZADO -----
\pagestyle{fancy}
\fancyhf{}
\lhead{\footnotesize Fundamentos de Mecánica (1000019) -- Taller}
\rhead{\footnotesize Semestre 2026-2}
\cfoot{\thepage}

% ----- LISTAS COMPACTAS -----
\setlist[itemize]{
	leftmargin=*,
	itemsep=2pt,
	topsep=2pt,
	parsep=0pt
}

\begin{document}
	
	% =========================================================
	% ENCABEZADO
	% =========================================================
	
	\begin{center}
		\vspace*{-0.8cm}
		\includegraphics[width=0.15\textwidth]
		{imagenes/universidad-nacional-de-colombia-sede-bogota-seeklogo.png}\\[0.1cm]
		
		{\Large\bfseries Taller de Movimiento Circular -- Grupo 14}\\[0.05cm]
		{\large\bfseries Fundamentos de Mecánica (1000019)}
	\end{center}
	
	
	% =========================================================
	% INFORMACIÓN DEL CURSO
	% ========================================================
	
	\vspace{0.2cm}
	
	% =========================================================
	\section{Movimiento Circular}
	% =========================================================
	
	\begin{enumerate}[label=\arabic*., leftmargin=*, itemsep=8pt]
		
		\item Dos puntos tienen coordenadas cartesianas $\left(1,\ \dfrac{\sqrt{3}}{2}\right)$ y $\left(-1,\ \dfrac{\sqrt{3}}{2}\right)$.
		\begin{enumerate}[label=\alph*)]
			\item Determine las coordenadas polares de esos dos puntos.
			\item Determine la distancia entre ellos.
		\end{enumerate}
		
		\item Tu compañero de cuarto está trabajando en su bicicleta y la pone de pies a cabeza. Enseguida notas que en el rin de $60$ cm de diámetro se encuentra atascada una piedra que gira a tres vueltas por segundo. ¿Cuáles son la rapidez y la aceleración de la piedra?
		
		\item Calcule la rapidez en km/h y la aceleración en m/s$^2$ de una persona sobre el ecuador, originadas por el movimiento rotacional de la Tierra. El radio de la Tierra es $r_T \simeq 6.37\times10^{6}$ m.
		
		\item En el modelo del átomo de hidrógeno propuesto por Niels Bohr, un electrón orbita alrededor de un protón estacionario en un movimiento circular uniforme, en un círculo de radio $r = 5.28\times10^{-11}$ m, con una rapidez $v = 2.18\times10^{6}$ m/s. Halle:
		\begin{enumerate}[label=\alph*)]
			\item La magnitud de la aceleración radial.
			\item El periodo del movimiento.
		\end{enumerate}
		
		\item El brazo de un aspersor de agua gira con una rapidez constante de tal forma que realiza $20$ vueltas en $5$ s. La distancia entre los extremos del brazo es de $20$ cm. Determine:
		\begin{enumerate}[label=\alph*)]
			\item El periodo del movimiento.
			\item La frecuencia angular.
			\item La rapidez del extremo del brazo.
			\item La aceleración centrípeta.
		\end{enumerate}
		Exprese sus resultados en unidades del SI y en unidades de $\pi$.
		
		\item La distancia entre Santo Tomé (África) y Quito es de aproximadamente $9500$ km. Ambas ciudades están casi sobre el ecuador. ¿Cuánto le toma a la Tierra rotar para cubrir la distancia entre las dos ciudades?
		
		\item Si la posición angular de un objeto en movimiento circular es $\theta(t) = \left(t^{2} + 2t + 1\right)$ rad, con $t$ en segundos, ¿cuáles son sus velocidad y aceleración angulares como función del tiempo?
		
	\end{enumerate}
	
	\vfill
	\noindent
	\begin{center}
		\small\textit{*Ejercicios tomados y adaptados del taller de la profesora Edna Carolina Pinilla Beltrán.}
	\end{center}
	
	\vspace{0.5cm}
	{\color{UnalGreen}\hrule height 0.5pt}
	\begin{center}
		\small
		\textit{Trabajo en grupos.}
	\end{center}
	
\end{document}
